\subsection{微分与导子}

设 K 为交换域。根据 III，第 558 页，命题 5， \(\mathbf{K}[(\mathbf{X}_i)_{i\in I}]\) 的每个导子 D 都以唯一方式延拓为导子 \(\bar{\mathbf{D}}\) （ \(\mathbf{K}((\mathbf{X}_i)_{i\in I})\) 的）。若 D, D' 是 \(\mathbf{K}[(\mathbf{X}_i)_{i\in I}]\) 的可交换导子，则括号 \([\mathbf{D},\mathbf{D}'] = \mathbf{DD}' - \mathbf{D}'\mathbf{D}\) 为零，因此 \([\bar{\mathbf{D}},\bar{\mathbf{D}}']\) （它是 \(\mathbf{K}((\mathbf{X}_i)_{i\in I})\) 的、延拓了 {[}D, D'{]} 的导子）为零；换言之， \(\bar{\mathbf{D}}\) 与 \(\bar{\mathbf{D}}'\) 是可交换的。特别地，导子 D\_i（IV，第 6 页）延拓为 \(\mathbf{K}((\mathbf{X}_i)_{i\in I})\) 的导子，仍记作 D\_i，且它们两两可交换。如果 \(f\in \mathbf{K}((\mathbf{X}_i)_{i\in I})\) , D\_if 也写作 D\_x\_i f 或 \(\frac{\partial f}{\partial\mathbf{X}_i}\) 或 \(f_{\mathbf{x}_i}'\) 。当只有一个未定元 X 时，使用记号 Df， \(\frac{df}{d\mathbf{X}}\) , \(f'\) 。

设 \(\mathbf{B} = \mathbf{K}[(\mathbf{X}_i)_{i \in I}]\) , \(\mathbf{C} = \mathbf{K}((\mathbf{X}_i)_{i \in I})\) 。根据 III，第 574 页，命题 23，典范映射

\[
\Omega_ {\mathrm{K}} (\mathrm{B}) \otimes_ {\mathrm{B}} \mathrm{C} \rightarrow \Omega_ {\mathrm{K}} (\mathrm{C})
\]

是向量 C-空间的一个同构。考虑到 III，第 570 页，我们看到向量 C-空间 \(\Omega_{\mathbf{K}}(\mathbf{C})\) 以族 \((d\mathbf{X}_i)_{i\in I}\) （由 \(\mathbf{X}_i\) 的微分组成）为一组基。设 \(\partial_i\) 是指标 \(i\) 的坐标形式（在 \(\Omega_{\mathbf{K}}(\mathbf{C})\) 上、相对于该基）。则映射 \(u\mapsto \langle \partial_i,du\rangle\) （C 到自身的）是 C 的一个导子，它把 \(\mathbf{X}_i\) 映到 1 ，把 \(\mathbf{X}_j\) 映到 0 （对 \(j\neq i\) ），因此等于 \(\mathrm{D}_i\) ；换言之，我们有

\[
d u = \sum_ {i \in I} \left(\mathrm{D} _ {i} u\right) d \mathrm{X} _ {i}\tag{1}
\]

对每个 \(u \in \mathbf{C}\) 成立。若 \(\mathbf{I}\) 是有限的，则 \((\mathrm{D}_i)_{i \in \mathbf{I}}\) 是 \(\mathbf{C}\) 的导子构成的向量 \(\mathbf{C}\) -空间的一组基。

命题 4. --- 设 E 是一个结合、交换且含单位元的 K-代数， \(\mathbf{x} = (x_i)_{i \in I}\) 是 E 的元素的一个族，且 \(f \in \mathbf{K}((\mathbf{X}_i)_{i \in I})\) 。假设 \(\mathbf{x}\) 在 \(f\) 中可代入，且 \(y = f(\mathbf{x})\) 。

\begin{enumerate}
\def\labelenumi{(\roman{enumi})}
\item
  对每个导子 \(\Delta\) （ \(\mathbf{K}((\mathbf{X}_i)_{i\in I})\) 的、把 \(\mathbf{K}[(\mathbf{X}_i)_{i\in I}]\) 映入自身的）， \(\mathbf{x}\) 在 \(\Delta f\) 中可代入。
\item
  对于E到某个E-模的每个导子D，我们有
\end{enumerate}

\[
\mathrm{D} y = \sum_ {i \in I} (\mathrm{D} _ {i} f) (\mathbf {x}) \cdot \mathrm{D} x _ {i}.
\]

设 \(f = \frac{u}{v}\) ，其中 \(u, v \in \mathbf{K}[(\mathbf{X}_i)_{i \in I}]\) 且 \(v(\mathbf{x})\) 在 E 中可逆。设 \(\Delta\) 是 \(\mathbf{K}((\mathbf{X}_i)_{i \in I})\) 的一个把 \(\mathbf{K}[(\mathbf{X}_i)_{i \in I}]\) 映入自身的导子。我们有

\[
\Delta f = \frac {(\Delta u) v - u (\Delta v)}{v ^ {2}}
\]

且 \(v^2(\mathbf{x})\) 可逆，因此 \(\mathbf{x}\) 在 \(\Delta f\) 中可代入。其次令 \(r = u(\mathbf{x})\) , \(s = v(\mathbf{x})\) ；我们有 \(y = s^{-1}r\) ，因此对 E 到某个 E-模的每个导子 D，我们有

\[
\begin{array}{l} \mathrm{D} y = s ^ {- 2} (s (\mathrm{D} r) - r (\mathrm{D} s)) \\ = s ^ {- 2} \left(s \sum_ {i \in I} (\mathrm{D} _ {i} u) (\mathbf {x}) \cdot \mathrm{D} x _ {i} - r \sum_ {i \in I} (\mathrm{D} _ {i} v) (\mathbf {x}) \cdot \mathrm{D} x _ {i}\right) \end{array}
\]

（据 IV，第 6 页，命题 4）。于是 \(Dy = \sum_{i \in I} w_i \cdot Dx_i\) ，其中

\[
w _ {i} = v (\mathbf {x}) ^ {- 2} (v (\mathbf {x}) \left(\mathrm{D} _ {i} u\right) (\mathbf {x}) - u (\mathbf {x}) \left(\mathrm{D} _ {i} v\right) (\mathbf {x})) = \left(\mathrm{D} _ {i} f\right) (\mathbf {x}).
\]

